% Generated draft: requires mathematical and source review.
\documentclass[11pt]{article}
\usepackage[margin=25mm]{geometry}
\usepackage{fontspec}
\setmainfont{Latin Modern Roman}
\usepackage{amsmath,amssymb,mathtools,unicode-math}
\setmathfont{Latin Modern Math}
\usepackage{enumitem,needspace,xurl,hyperref}
\setlength{\parindent}{0pt}
\setlength{\parskip}{5pt}
\setlength{\emergencystretch}{3em}
\newcommand{\N}{\mathbb{N}}
\newcommand{\Z}{\mathbb{Z}}
\newcommand{\Q}{\mathbb{Q}}
\newcommand{\R}{\mathbb{R}}
\newcommand{\C}{\mathbb{C}}
\newcommand{\F}{\mathbb{F}}
\newcommand{\A}{\mathbb{A}}
\newcommand{\PP}{\mathbb{P}}
\newcommand{\E}{\mathbb{E}}
\newcommand{\eps}{\varepsilon}
\newcommand{\dd}{\,\mathrm{d}}
\begin{document}
\section*{UnsolvedMath ID: 30006025}
\textbf{Source problem code:} catalogue entry
\subsection{Definitions and mathematical statement}
For terminology, an embedded graph \(\Gamma\subset S\) is called a \emph{one-face map} when its complement \(S\setminus\Gamma\) is a single topological open disk. A finite tree is a finite connected graph without cycles; a \emph{decorated tree} is a tree equipped with additional data. The supplied record does not specify the allowed surfaces, graph conventions, rooting conventions, decorations, or equivalence relations for the maps and trees in the intended version of Chapuy's bijection.

\textbf{Construction problem.} Find a geometric adaptation of the bijection attributed to Chapuy between one-face maps and decorated trees that constructs random surfaces. A precise formulation requires identification of the particular bijection, its domain and codomain, the geometric structures on the resulting surfaces, the probability distributions involved, and the sense in which the construction realizes the combinatorial correspondence. These specifications are absent from the supplied record, so the construction request cannot presently be expressed as a uniquely determined quantified mathematical statement.

\textbf{Related enumerative question.} Find a geometric explanation of the enumerative fact referred to in the original excerpt as being key to Theorems 1 and 2. The excerpt ends immediately before that fact is stated, and neither theorem is supplied. Consequently, no particular enumerative identity can be stated here. The original excerpt presents this as a separate, more modest question; the cleaned statement combines it with the geometric construction request, but does not supply the missing identity.
\subsection{Short English statement}
Seek a geometric version of Chapuy's correspondence between one-face maps and decorated trees that constructs random surfaces. A separate question asks for a geometric explanation of an enumerative fact omitted from the supplied excerpt. The intended map classes, decorations, surface geometry, probability laws, and enumerative fact must be recovered before either request has a precise mathematical formulation.
\subsection{Sources}
\begin{itemize}
\item Supplied local UnsolvedMath registry record, permanent ID \(30006025\), title \emph{Geometric Chapuy Bijections for Random Surfaces}, source problem code \texttt{catalogue entry}, in the Oberwolfach Reports Open Problems collection. The record attributes the combinatorial bijection to Chapuy; no problem proposer is identified.
\item Supplied report citation: \emph{Riemann Surfaces: Random, Flat, and Hyperbolic Geometry} (2025), Oberwolfach Reports, pp.\ 2369--2424, DOI \texttt{10.4171/owr/2024/41}. Source URL and supplied literature source labeled ``OWR report'': \url{https://doi.org/10.4171/owr/2024/41}. This is the supplied report-level reference; no exact problem page or complete bibliography entries for the excerpt's references [4] and [6] are provided.
\item The original excerpt identifies Chapuy's bijection as reference [4], and says it is a key tool in the proof of reference [6]. It also refers to Theorems 1 and 2 and a subsequent enumerative fact. Those references, theorem statements, and the enumerative fact are not included in the supplied record.
\item Literature assessment checked at \texttt{2026-08-22}: the supplied curation labels the problem open, states that the geometric adaptation is explicitly left open, and reports that no primary construction resolving it was located. This is supplied dated curation, not an independently verified literature assessment. The background explicitly requests independent verification of time-sensitive status and cited claims.
\item The registry labels its statement \texttt{corrected\_verified} and reports that the official report was checked and the question isolated and its notation repaired. These are supplied verification claims, not verification performed for this draft. The record gives proposed year 2024, while its report citation displays 2025; these supplied year designations are retained without reconciliation.
\end{itemize}
\end{document}
